% !TEX root = master.tex
% !TEX program = lualatex

\section{Introduction}
(L1 + L2)

\subsection{Goal of the Course}

The goal of this course is to understand how it is possible that when a user clicks in a web browser, a resource located somewhere else on the Internet appears on the user's screen.

Example scenario:
\begin{itemize}
\item A user clicks on a video in a browser.
\item The browser communicates with a remote server.
\item The server sends the requested data.
\item The browser displays the video.
\end{itemize}

This apparently simple action actually involves many components of a computer system and the Internet.

\subsection{The Main Components of the Internet}

The Internet is composed of three fundamental elements:

\begin{itemize}
\item \textbf{End-systems}
\item \textbf{Packet switches}
\item \textbf{Network links}
\end{itemize}

\subsubsection{End-systems}

An \textbf{end-system} is a device connected to the Internet.

Examples include:

\begin{itemize}
\item laptops
\item smartphones
\item servers
\item connected appliances
\item cars
\item medical devices
\end{itemize}

End-systems run applications and communicate with other end-systems.

\subsubsection{Packet switches}

A \textbf{packet switch} is a device that helps connect end-systems together and forward packets across the network.

Examples:
\begin{itemize}
\item routers
\item switches
\end{itemize}

Packet switches route packets toward their destination.

\subsubsection{Network links}

A \textbf{network link} is the physical medium that connects two devices.

Examples include:

\begin{itemize}
\item fiber optic cables
\item Ethernet cables
\item WiFi radio signals
\end{itemize}

\subsection{Programs, Processes and Threads}

\subsubsection{Programs}

A \textbf{program} is a set of files stored on disk.

Examples:
\begin{itemize}
\item a web browser
\item a ping utility
\item a git client
\item a messaging application
\end{itemize}

\subsubsection{Processes}

When a program is executed, the operating system creates a \textbf{process}.

Definition:

\begin{quote}
A process is an instance of a program executing in memory.
\end{quote}

Processes live in main memory.

Multiple processes can correspond to the same program.

\subsubsection{Threads}

A process can contain one or more \textbf{threads}.

A thread is a unit of execution within a process.

Analogy:

\begin{itemize}
\item Program = musical score
\item Process = orchestra playing the score
\item Threads = individual instruments
\end{itemize}

\subsection{Distributed Applications}

A \textbf{distributed application} is composed of:

\begin{itemize}
\item multiple processes
\item running on potentially different computers
\item communicating through messages
\item working toward a common goal
\end{itemize}

Example:

\begin{itemize}
\item browser process
\item YouTube server process
\end{itemize}

These processes exchange messages over the Internet.

\subsection{Communication Protocols}

A \textbf{communication protocol} defines the rules governing interactions between processes.

Definition:

\begin{quote}
A communication protocol specifies all possible message sequences that can be exchanged between communicating parties.
\end{quote}

Example protocol:

\begin{enumerate}
\item Client: hello
\item Server: hello back
\item Client: request file
\item Server: send file
\end{enumerate}

Protocols provide:

\begin{itemize}
\item an interface between processes
\item an abstraction of one process to another
\end{itemize}

%%%%%%%%%%%%%%%%%%%%%%%%%%%%%%%%%%%%%%%%%%%%%%%%%%%%%%
\subsection{System Calls}

Processes interact with external resources through \textbf{system calls} (syscalls).

Examples:

\begin{itemize}
\item send()
\item recv()
\item read()
\item write()
\end{itemize}

System calls allow processes to access:

\begin{itemize}
\item network
\item storage
\item devices
\end{itemize}

Processes do not directly control these resources.

Instead, they request access through the operating system.

%%%%%%%%%%%%%%%%%%%%%%%%%%%%%%%%%%%%%%%%%%%%%%%%%%%%%%
\subsection{Operating System}

The \textbf{Operating System (OS)} is a special program that enables other programs to run.

Characteristics:

\begin{itemize}
\item loaded in memory during boot
\item intermediary between processes and hardware
\item reacts to system calls
\item reacts to hardware events
\end{itemize}

The OS coordinates:

\begin{itemize}
\item CPU
\item memory
\item network
\item storage
\end{itemize}

%%%%%%%%%%%%%%%%%%%%%%%%%%%%%%%%%%%%%%%%%%%%%%%%%%%%%%
\subsection{The I/O Stack}

Process interaction with devices occurs through layers.

Example for storage:

\begin{itemize}
\item Process
\item File System
\item Device driver
\item Device
\end{itemize}

Example for networking:

\begin{itemize}
\item Process
\item Network subsystem
\item Device driver
\item Network device
\end{itemize}

This layered structure is called the \textbf{I/O stack}.

Advantages:

\begin{itemize}
\item abstraction
\item modularity
\item flexibility
\item separation of concerns
\end{itemize}

